\section{Discussion}\label{sec:discussion}

This section discusses lessons learned from our experiments and past experiences. We then present several ideas for future improvements. 

\subsection{Lessons Learned}

Distributed data parallel training is a conceptually simple or practically subtle framework. There are various techniques to improve its speed, creating a complex configuration space. Based on our observations, there is no single configuration that would work for all use cases, as it would highly depend on the model size, model structure, network link bandwidth, etc. However, on individual configuration dimensions, we summarized intuitions to help application developers to quickly navigate to a small range which likely contains the optimal solution for a given use case. The specific value of the configuration would require empirical measurements for every different deployment. 

\begin{itemize}
    \item \textbf{Communication Backend}: NCCL is considerably faster than Gloo in most use cases. When available, applications should seek to use NCCL as the primary collective communication backend.
    \item \textbf{Bucket Size}: Both excessively small or large bucket sizes are detrimental to communication performance. The optimal value lives in between and depends on the type of communication backend employed. The optimal bucket sizes are likely to increase with the size of the model in a sub-linear manner. 
    \item \textbf{Resource Allocation}: There is a significant slowdown with NCCL backend when scaling models across machine boundaries, if the bandwidth across machines is considerably lower than that between same-machine GPUs. In such cases, it is recommended to keep the \code{DDP} group within the same machine. If the training requires larger scale, developers can explore enabling \code{no\_sync} mode if it attains acceptable convergence speed. 
\end{itemize}

\subsection{Future Improvements}

While we implement and maintain the \code{DDP} package, several ideas for improvements popped up. This section discusses the basic ideas behind those improvements. 

%\subsubsection{Heterogeneous Bucket Size}

%Current \code{DDP} implementation uses the same size for all buckets. It simplifies the configuration for coalescing communications but might lead to unnecessary network idling. In practice, some applications set the bucket size to a few hundred \code{MB}s to achieve the highest communication efficiency. However, as a result, the first few \code{AllReduce} operations only launches after a several hundred \code{MB}s of gradients are ready, which could take the autograd engine a while to reach that point. Moreover, as gradient bucketing order does not faithfully represent the exact backward computation order, \code{DDP} might need to wait even longer to start communicating.  To alleviate this problem, \code{DDP} can employ heterogeneous bucket sizes, where the first few buckets can be smaller than the configured size cap. In this way, communication can start earlier and can overlap more with computations.

\subsubsection{Gradient Order Prediction}

Although \code{DDP} cannot deterministically detect the backward computation order on all parameters at construction time, the order usually does not change that often in practice. One viable solution is to trace the backward order using autograd hooks and update parameter to bucket mapping accordingly. As bucket re-allocation will introduce noticeable overhead, it should be conducted infrequently. Given the existing complexities in \code{DDP}, tracing overhead should be negligible. Nevertheless, if there are disparities among tracing results from different iterations, additional complexities will be necessary to reach a consensus. 

\subsubsection{Layer Dropping}

One technique to accelerate training and avoid overfitting is to randomly drop layers during the forward pass~\cite{fan2019reducing}. This works well with local training. As every forward pass would build a new autograd graph, those skipped layers will not participate in the backward pass either. This idea also works with \code{DDP}, because parameters in skipped layers can be marked as ready in the forward pass and \code{DDP} will not wait for their autograd hooks during the backward pass. Although \code{DDP} would produce the correct result, this technique alone is inadequate to accelerate distributed data parallel training the same way as local training due to the fixed parameter-to-bucket mapping. As \code{AllReduce} uses a bucket as the minimum granularity, it cannot judiciously react to vacancies in buckets (\emph{i.e.}, skipped layers or parameters). Consequently, regardless of how the forward pass skips layers, there is always the same amount of data to be communicated across the wire during the backward pass. Besides, \code{DDP} cannot afford the luxury to adjust all buckets to cooperate with randomly skipped layers, as that would result in unacceptable memory allocation overhead. To tackle this problem, one solution is to keep bucket buffers intact but modify parameter-to-bucket mappings accordingly. Another option is to perform layer skips at the bucket level, \emph{i.e.}, \code{DDP} can map layers instead of parameters to buckets and all processes skip the same bucket in the same iteration. Both options require extra coordination across all \code{DDP} processes, which can be implemented by using the same random seed or having an authority process to broadcast the plan.

\subsubsection{Gradient Compression}

Another potential improvement for \code{DDP} is to reduce the volume of data for communication by compressing gradients. The absolute value of gradients are usually small, which might not require \code{float32} or \code{float64} types. Current \code{DDP} implementation always uses the parameter type as the gradient type that can become an overkill especially when the model is approaching convergence. In this case, \code{DDP} would benefit from adaptive compression levels by only communicating gradients with the necessary precision. Some recent research work~\cite{seide2014} even proposes more aggressive compression schemes, where by trading a tiny amount of model accuracy, applications can significantly accelerate distributed training by communicating just 1 bit for each gradient. 

%\subsubsection{Extensible Design}

%Existing \code{DDP} implementations are closely coupled in both \code{distributed.py} and \code{reducer.cpp}. Any addition or modification to the algorithm would require a thorough understanding of \code{DDP} code. Furthermore, many production and research projects have integrated \code{DDP} into their critical code paths which are sensitive to any performance regressions. The current structure allows us to quickly iterate on the initial implementation and to enable various applications in vision, NLP, and recommendation systems, but the design and dependencies gradually drive \code{DDP} to be a rigid and monolithic chunk of code that is reluctant to embrace new features. For example, even if we know that the improvements discussed the in previous sections can certainly help different use cases, they will not easily land into \code{DDP} as they would slow down some other training scenarios or pollute the \code{DDP} API. It has become a daunting hurdle for the open-source community to try out and contribute ideas and implementations as well. To open up opportunities for future improvements and customization, \code{DDP} requires an overhaul on the structure to dismantle and rebuild in a modularized and pluggable way.